% The research corpus (ci/research.json): a humanities article in
% KOMA-Script (scrartcl) with babel, csquotes and biblatex's authoryear
% style through biber, on biblatex's own example database.
\documentclass[a4paper,11pt,parskip=half,bibliography=totoc]{scrartcl}
\usepackage[T1]{fontenc}
\usepackage[ngerman,main=english]{babel}
\usepackage{lmodern}
\usepackage{csquotes}
\usepackage[style=authoryear-comp,backend=biber,maxcitenames=2]{biblatex}
\addbibresource{biblatex-examples.bib}
\usepackage{scrlayer-scrpage}
\pagestyle{scrheadings}
\ihead{Reading Aristotle}
\ohead{\pagemark}
\cfoot{}
\title{Reading Aristotle Today}
\subtitle{Notes on Reception and Translation}
\author{H. Humanist}
\date{1 January 2026}
\begin{document}
\maketitle
\begin{abstract}
A short essay in the shape of a humanities paper: author--year
citations, footnotes, block quotations and a German passage.
\end{abstract}
\section{Introduction}
Aristotle's \textit{Poetics} has been read in many ways
\autocite{aristotle:poetics}; on the soul he is read differently again
\autocites[12]{aristotle:anima}[see also][]{aristotle:physics}.\footnote{On
the translations see \textcite{aristotle:rhetoric}, and for a modern
reading \textcite{nussbaum}.}
\section{Quotation}
\blockquote[{\cite[45]{augustine}}]{Every long quotation, set off from the
text, is introduced by its author and closed by its source, as the
humanities ask; it runs on for several lines so that the block
quotation's margins can be seen at work in the finished page.}
\section{A German passage}
\begin{otherlanguage}{ngerman}
Die \enquote{Übersetzung} ist, wie \textcite{nietzsche:ksa1} zeigt, nie nur
Übertragung; Silbentrennung und Anführungszeichen folgen dem Deutschen.
\end{otherlanguage}
\section{Conclusion}
Further reading: \cite{kant:kpv,kant:ku,worman}.
\printbibliography
\end{document}
